\chapter{Environnement technique}
\label{chap:env}

% =====================================================================
%  RAPPEL INSTRUCTION : présenter SUCCINCTEMENT les technologies, outils,
%  méthodes. Éviter l'inventaire lassant : ne lister que ce qui est
%  réellement mobilisé dans les tâches, et regrouper par thème.
% =====================================================================

\section{Langages et bibliothèques}
% TODO : langages utilisés (ex. Python/C/C++/Shell...), frameworks de test
% (ex. pytest, ...), bibliothèques clés. Une phrase de justification par item.
\emph{[À rédiger.]}

\section{Outils de build et de compilation}
% TODO : chaîne de compilation, SDK, gestionnaires de dépendances,
% (utile en amont de la tâche 3 « rétro-ingénierie »).
\emph{[À rédiger.]}

\section{Gestion de version et intégration continue}
% TODO : \gls{vcs} (Git ?), forge (GitLab/GitHub ?), pipelines \gls{ci}/\gls{cd}.
\emph{[À rédiger.]}

\section{Méthodologie et gestion de projet}
% TODO : méthode (agile/scrum/kanban ?), outils de ticketing, rituels.
\emph{[À rédiger.]}
